% =====================================================================
% Ensaio 1 — Seção 3: Referencial teórico
%
% Vai do geral ao específico e TERMINA numa predição testável — que é o que
% amarra esta seção à §4 (identificação) e planta o Ensaio 2.
%
% ⚠️ REFERÊNCIAS NÃO VERIFICADAS. As citações abaixo vêm de
% docs/framework-dissertacao.md, documento do pesquisador. Nenhuma foi conferida
% contra DOI/Crossref — as bases estão inacessíveis (ver docs/lacunas-de-dados.md
% §1, classe E). Conferir antes da versão final; o CLAUDE.md é explícito.
% =====================================================================

\section{Referencial teórico}
\label{sec:teoria}

\subsection{A deriva como externalidade negativa}
\label{sub:externalidade}

O agrotóxico que deriva sobre o vizinho é o caso de manual de externalidade
negativa: um custo imposto a terceiros que não transita pelo sistema de preços. O
produtor que contrata a aplicação aérea internaliza o custo do serviço e do
produto, não o custo em saúde suportado pela população a favor do vento. O
resultado privado é aplicação acima do nível socialmente eficiente.

As duas respostas canônicas divergem no instrumento e, mais fundo, no que
consideram o objeto do problema. A resposta pigouviana internaliza o custo por
tributo igual ao dano marginal, fazendo o mercado precificar o que antes não
precificava. A resposta coaseana observa que, com direitos de propriedade bem
definidos e custos de transação baixos, as partes negociam a solução eficiente
por conta própria --- e desloca a pergunta para \emph{quem detém o direito}: o de
pulverizar, ou o de não ser pulverizado.

\paragraph{Por que a barganha coaseana não se aplica aqui.} A deriva atinge
população rural difusa, numerosa e não organizada, sobre território extenso, com
dano probabilístico, latente e de difícil atribuição individual. Os custos de
transação são, por construção, proibitivos. Isso não é detalhe de aplicação: é o
que \emph{justifica} a intervenção regulatória em vez da negociação privada, e é
também o que desloca a discussão para a escolha de instrumento --- objeto do
Ensaio 2.

\subsection{Preço ou quantidade}
\label{sub:weitzman}

O núcleo teórico do Ensaio 2 é a comparação de instrumentos sob incerteza. Sob
incerteza a respeito dos custos de controle, um instrumento de preço --- taxa
sobre o insumo --- e um de quantidade --- cota, ou proibição --- \textbf{não são
equivalentes em bem-estar esperado}. Qual deles domina depende das
\textbf{inclinações relativas} das curvas de benefício e custo marginais.

A regra de bolso: se o benefício marginal do controle é mais inclinado que o
custo marginal --- danos que disparam rapidamente ao ultrapassar um limiar,
plausível em desfechos de saúde ---, o instrumento de \textbf{quantidade} domina.
Se o custo marginal é mais inclinado, o instrumento de \textbf{preço} domina. Uma
proibição é o caso-limite de cota fixada em zero.

\begin{quote}
\textbf{A ligação que amarra os dois ensaios.} Weitzman decide pela
\textbf{inclinação}, não pelo nível. É por isso que o Ensaio 1 tem como
parâmetro-alvo a \emph{curva} dose-resposta e não o efeito num ponto: o Ensaio 2
consome a derivada. Um Ensaio 1 que entregasse apenas o efeito médio no nível
tratado deixaria o Ensaio 2 sem âncora empírica --- e a escolha entre proibir e
taxar voltaria a ser argumentada por plausibilidade, que é exatamente o que a
dissertação se propõe a substituir.
\end{quote}

\paragraph{A camada de realismo: poluição difusa.} Instrumentos baseados em
emissão observada pressupõem que se sabe quem emitiu quanto. A deriva não tem
emissor identificável --- é poluição de fonte difusa. Sob heterogeneidade entre
produtores, a literatura registra que a taxa tende a ser mais eficiente enquanto
a regulação por quantidade gera maior lucro ao setor regulado --- resultado com
consequências distributivas, e não apenas de eficiência. A observação importa
porque a proibição cearense foi contestada judicialmente pelo lobby do setor, e
não por consumidores.

\subsection{A regulação de agrotóxicos em específico}
\label{sub:regulacao}

A literatura específica sobre regulação de pesticidas acrescenta duas coisas que
o instrumental geral não dá. A primeira é a análise marginal dos custos de
bem-estar \emph{e dos efeitos distributivos} da regulação --- quem paga a conta
da restrição, e em que proporção. A segunda é a comparação empírica de
instrumentos, em que há registro de cenários nos quais cotas superam taxas ---
isto é, o resultado de Weitzman aplicado, e não apenas invocado.

\paragraph{Uma advertência que a §\ref{sub:quimico} torna necessária.} A
literatura econômica sobre regulação de agrotóxicos é, majoritariamente,
construída sobre herbicidas em culturas de grão, frequentemente transgênicas.
O caso cearense é de fungicidas e inseticidas sobre fruticultura irrigada. A
estrutura de custos de substituição, a elasticidade de resposta do produtor e o
perfil toxicológico são outros. As predições qualitativas do referencial se
mantêm; as magnitudes emprestadas, não.

\subsection{O que o instrumento pressupõe, conta e exclui}
\label{sub:normativo}

Escolher entre proibir, taxar e regular desempenho é problema de desenho de
instrumento. Tratá-lo apenas como exercício de maximização de bem-estar esperado,
porém, deixa fora do quadro uma camada que este caso torna difícil de ignorar.

Precificar uma externalidade de saúde --- a via da taxa --- e proibi-la --- a via
da cota zero --- não são escolhas técnicas equivalentes com sinais trocados.
Elas diferem no que pressupõem: a primeira supõe que dano à saúde é comensurável
com dinheiro e que a compensação é a forma adequada de resposta; a segunda supõe
que há um direito a não ser exposto, anterior ao cálculo. Diferem também no que
contam: um cálculo welfarista puro registra o custo agregado da restrição e o
benefício agregado em saúde, e não registra quem, especificamente, absorve o
resíduo quando o instrumento falha.

A lei em análise leva o nome de um líder rural assassinado por se opor à
pulverização. Uma avaliação que se limite a estimar o efeito médio sobre peso ao
nascer não está errada --- é, aliás, o que esta dissertação faz, e faz com o rigor
que o método exige. Mas ela é incompleta se não registrar que a escolha de
instrumento é também uma posição sobre de quem é o direito. Esta seção declara a
camada; a §\ref{sec:identificacao} não a usa como argumento empírico, e a
conclusão retorna a ela na discussão de política.

\subsection{A predição testável}
\label{sub:predicao}

O referencial acima produz predições, e é por elas que a §\ref{sec:identificacao}
é avaliada.

\begin{description}
  \item[Predição principal.] Se a pulverização aérea gera externalidade de saúde
    por deriva, sua proibição reduz a exposição populacional, e essa redução é
    \textbf{crescente na intensidade pré-ban} de aplicação. Municípios que mais
    aplicavam devem apresentar a maior melhora em desfechos perinatais. É a
    predição que a curva dose-resposta testa.

  \item[Predição de canal --- água.] Se parte da exposição transita pelo canal
    hídrico, a melhora deve ser maior a jusante do que a montante dentro da mesma
    bacia, para dose pré-ban comparável.

  \item[Predição de canal --- ar.] Se parte transita pelo canal atmosférico, a
    melhora deve ser maior a favor do vento dominante do que contra.

  \item[Predição de substituição, de sinal oposto.] Se a aplicação migra do
    avião para o solo, a intoxicação aguda \textbf{ocupacional} deve
    \textbf{aumentar} com a dose pré-ban, ainda que a exposição populacional
    caia.

  \item[Predição nula, que funciona como falsificação.] Como a proibição alcança
    o \emph{método} e não a \emph{molécula}, a disponibilidade do produto não se
    altera. Logo, a intoxicação \textbf{autoprovocada} --- que responde à
    disponibilidade, não ao modo de aplicação --- \textbf{não} deve se mover com
    a dose. As duas séries partilham população, sistema de registro e municípios;
    movimento conjunto não sustenta a hipótese de substituição, e sim algum fator
    que desloca internação em geral.

  \item[Predição de custo.] A proibição impõe custo ao produtor --- substituição
    de tecnologia de aplicação, possível perda de produtividade. É o lado do
    trade-off que o Ensaio 2 precifica, e cuja mensuração fica na agenda.
\end{description}

\paragraph{O que torna o conjunto informativo.} As predições não apontam todas na
mesma direção. Duas delas --- substituição e falsificação --- têm sinais
esperados que \emph{contradizem} a leitura ingênua de que "a proibição melhora
tudo". Um desenho no qual todo resultado possível confirmaria a hipótese não
testaria hipótese alguma; a §\ref{sub:poder} formaliza esse critério ao declarar,
antes da estimação, que resultado constituiria evidência contrária.
